% Section 2: Multilinear maps} | Part: 1 | Depends on: see docs/REWRITE_PLAN.md | Status: todo | Original lines: 511-545
%!TEX root = ../../../main.tex
\subsection{Multilinear maps}\label{sec:multilinear}
\begin{formal}[{multilinear map}]
Let $E_1,\dots,E_k$ be vector spaces. A map $f:E_1\times\dots\times E_k\to\R$ is \emph{multilinear} if it is linear in each slot when the other slots are held fixed:
\[
f(v_1,\dots,u+t\tilde u,\dots,v_k)=f(v_1,\dots,u,\dots,v_k)+t\,f(v_1,\dots,\tilde u,\dots,v_k).
\]
Examples: a linear form ($k=1$) is an element of $E^*$; the dot product $\langle v_1,v_2\rangle$ is bilinear ($k=2$); the determinant of $m$ column vectors is $m$-linear.
\end{formal}

\begin{intuition}[{linear in every slot = ``a machine with several input slots''}]
A multilinear map is a machine with $k$ input hoppers. You may feed any vector in each hopper; the output is a number that depends linearly on what is in each hopper \emph{separately}. Multilinear maps are exactly the maps that are completely determined by their values on basis vectors.
\end{intuition}

\begin{definition}[Multilinear map]
Let $V_1,\dots,V_k,U$ be vector spaces. A map $\beta:V_1\times\cdots\times V_k\to U$ is
\emph{multilinear} if it is linear in each argument when the others are held fixed.
For $k=2$ we say \emph{bilinear}. If all $V_i=V$ and
$\beta(\dots,v_i,\dots,v_j,\dots)=-\beta(\dots,v_j,\dots,v_i,\dots)$ for all $i\neq j$,
$\beta$ is \emph{alternating}; we write $\Alt^k(V;U)$ for the space of such maps.
\end{definition}

\begin{example}
The determinant $\det:(\K^n)^n\to\K$ (as a function of the $n$ columns) is multilinear and alternating.
The dot product $\R^n\times\R^n\to\R$ is bilinear (and symmetric).
\end{example}

The crucial remark is that a bilinear map is \emph{not} a linear map on $V\times W$:
$\beta(2u,2w)=4\beta(u,w)$, whereas a linear map would give $2\beta(u,w)$. The tensor product
will fix exactly this defect.
